% Capítulo 1 - Introdução

\chapter{Introdução}
\label{cap:introducao}

A quantidade de dados produzidos em ambientes domésticos cresce continuamente. Computadores pessoais, celulares, câmeras, aplicativos de mensagens, consoles, bibliotecas 
multimídia e rotinas de backup geram um volume cada vez maior de fotografias, vídeos, documentos e arquivos pessoais. Nesse cenário, o armazenamento deixa de ser apenas 
um recurso auxiliar e passa a ocupar papel estrutural na vida digital cotidiana.

Serviços de nuvem pública oferecem conveniência, sincronização automática e acesso remoto, mas também criam dependência da conectividade externa, da política comercial do fornecedor e da disponibilidade da infraestrutura de terceiros. O objetivo deste trabalho não é rejeitar a nuvem pública, mas investigar uma alternativa local complementar para armazenamento e sincronização de dados.

O projeto propõe a construção e análise de um servidor NAS doméstico baseado em uma Radxa ROCK 4B+, três SSDs conectados por adaptador externo USB-SATA, OpenMediaVault, Nextcloud e Tailscale. O OpenMediaVault será utilizado como camada de administração do armazenamento em rede; o Nextcloud fornecerá recursos de nuvem privada; e o Tailscale permitirá acesso remoto seguro por meio de uma rede privada autenticada.

\begin{figure}[htb]
    \legenda[fig:arquitetura]{Arquitetura lógica da solução proposta}
    \fig{width=0.9\textwidth}{assets/img/arquitetura}
    \fonteautoria
\end{figure}

\section{Problema de Pesquisa}

É tecnicamente viável implementar um servidor NAS doméstico de baixo custo baseado em Single Board Computer, administrado por OpenMediaVault, complementado por Nextcloud e acessado remotamente via Tailscale, mantendo desempenho, acessibilidade e usabilidade adequados ao contexto doméstico?

\section{Hipótese}

Um servidor NAS doméstico baseado em SBC, com OpenMediaVault como camada de administração, Nextcloud como camada de nuvem privada e Tailscale como mecanismo de acesso remoto seguro, pode atender de forma satisfatória cenários domésticos de armazenamento, sincronização e compartilhamento, com custo inferior ao de NAS comerciais de entrada.

\section{Objetivo Geral}

Analisar a viabilidade técnica de um servidor NAS doméstico de baixo custo baseado em Single Board Computer, OpenMediaVault, Nextcloud e acesso remoto seguro.

\section{Objetivos Específicos}

\begin{alineas}
    \item Montar a infraestrutura física utilizando Radxa ROCK 4B+, fonte adequada, rede cabeada e três SSDs de 2,5 polegadas conectados por adaptador externo USB-SATA;
    \item Instalar e configurar Armbian/Debian minimal compatível com a plataforma escolhida;
    \item Instalar e configurar o OpenMediaVault como camada central de gerenciamento do NAS;
    \item Integrar os três SSDs, padronizando nomenclatura, montagem e organização lógica;
    \item Implantar compartilhamento local por SMB/CIFS para acesso em rede cabeada;
    \item Implantar o Nextcloud como camada de nuvem privada para sincronização e acesso web;
    \item Implantar o Tailscale via Compose para acesso remoto seguro ao ambiente doméstico;
    \item Executar testes locais e remotos de desempenho, estabilidade e usabilidade;
    \item Comparar o custo estimado da solução com modelos comerciais de NAS doméstico.
\end{alineas}

\section{Justificativa}

A justificativa acadêmica do trabalho está na combinação entre baixo custo de entrada, uso de software livre amplamente documentado, possibilidade de medição objetiva de desempenho e relevância prática para usuários domésticos. Diferentemente de um tutorial, o estudo pretende produzir evidências mensuráveis sobre desempenho de rede, comportamento do armazenamento, acesso remoto, consumo de recursos e comparação econômica.

\begin{figure}[htb]
    \legenda[fig:fluxos]{Fluxos de acesso local, remoto e administrativo}
    \fig{width=0.9\textwidth}{assets/img/fluxos}
    \fonteautoria
\end{figure}